\chapter{Ứng dụng đạo hàm để khảo sát và vẽ đồ thị hàm số}

\section{Tính đơn điệu và cực trị của hàm số}

\exercise Cho hàm số $y = x^3 - 3x^2 - 9x + 5.$
\begin{enumerate}
   \item Xét tính đơn điệu (tìm các khoảng đồng biến, nghịch biến) của hàm số.
   \item Tìm các điểm cực trị và giá trị cực trị tương ứng của hàm số.
\end{enumerate}

\exercise Tìm các khoảng đơn điệu và cực trị của hàm số phân thức
$y = \frac{x^2 - 2x + 4}{x - 2}.$

\exercise Tìm tất cả các giá trị thực của tham số $m$ để hàm số
$y = \frac{1}{3}x^3 - mx^2 + (m+2)x - 1$
đồng biến trên $\mathbb{R}$.

\section{Giá trị lớn nhất và giá trị nhỏ nhất của hàm số}

\exercise Tìm giá trị lớn nhất và giá trị nhỏ nhất của hàm số
$y = 2x^3 - 3x^2 - 12x + 1$
trên đoạn $[-2; 3]$.

\exercise Tìm giá trị nhỏ nhất của hàm số
$y = x + 1 + \frac{4}{x - 1}$
trên khoảng $(1; +\infty)$.

\exercise Người ta muốn làm một bể cá không nắp dạng hình hộp chữ nhật có đáy là hình vuông và thể tích bằng \qty{108}{dm^3}. Hãy xác định kích thước của bể cá (cạnh đáy và chiều cao) sao cho diện tích kính dùng để làm bể cá là nhỏ nhất (bỏ qua độ dày của kính).

\section{Đường tiệm cận của đồ thị hàm số}

\exercise Tìm các đường tiệm cận đứng và tiệm cận ngang của đồ thị hàm số
$$y = \frac{3x - 1}{x + 2}.$$ 

\exercise Tìm các đường tiệm cận của đồ thị hàm số
$$y = \frac{x^2 - 3x + 3}{x - 1}.$$ 

\exercise Một nhà sản xuất ước tính tổng chi phí để sản xuất $x$ sản phẩm là
$C(x) = 3x + 120\quad \text{(triệu đồng)}.$
Chi phí sản xuất trung bình cho mỗi sản phẩm là
$f(x) = \frac{C(x)}{x}.$ 
\begin{enumerate}
   \item Tìm đường tiệm cận ngang của đồ thị hàm số $y = f(x)$.
   \item Nêu ý nghĩa của đường tiệm cận ngang tìm được khi số lượng sản phẩm $x$ ngày càng lớn ($x \to +\infty$).
\end{enumerate}

\section{Khảo sát sự biến thiên và vẽ đồ thị của hàm số}

\exercise Khảo sát sự biến thiên và vẽ đồ thị của hàm số
$$y = -x^3 + 3x^2 - 2.$$

\exercise Khảo sát sự biến thiên và vẽ đồ thị của hàm số
$$y = \frac{x^2 - x - 1}{x - 2}.$$ 

\section{Ứng dụng đạo hàm giải quyết bài toán thực tiễn}

\exercise Số lượng vi khuẩn trong một thí nghiệm sinh học sau $t$ giờ nuôi cấy được mô hình hóa bởi hàm số
$$P(t) = \frac{10000}{1 + 9e^{-0,5t}} \quad (t \ge 0).$$
\begin{enumerate}
   \item Tính số lượng vi khuẩn ban đầu lúc bắt đầu nuôi cấy ($t = 0$).
   \item Tốc độ tăng trưởng vi khuẩn tại thời điểm $t$ được xác định bởi đạo hàm $P'(t)$. Sau bao nhiêu giờ kể từ lúc nuôi cấy thì tốc độ tăng trưởng vi khuẩn đạt giá trị lớn nhất?
\end{enumerate}

\exercise Một sân vận động có sức chứa \num{40000} khán giả. Khi giá vé thi đấu là $100$ nghìn đồng, số khán giả trung bình mua vé là \num{20000} người. Qua thăm dò thị trường, ban tổ chức nhận thấy nếu giá vé giảm thêm $10$ nghìn đồng thì số khán giả mua vé sẽ tăng thêm \num{2500} người. Giả sử số khán giả mua vé $x$ phụ thuộc tuyến tính vào giá vé $p$ (nghìn đồng).
\begin{enumerate}
   \item Lập biểu thức hàm số doanh thu bán vé $R(p)$ theo giá vé $p$.
   \item Ban tổ chức nên quy định giá vé là bao nhiêu để doanh thu từ bán vé đạt giá trị lớn nhất? Doanh thu cực đại thu được là bao nhiêu?
\end{enumerate}

